\section{变量分离方程与变量变换}

\subsection{变量分离方程}

\begin{mydefinition}{定义}
形如
\begin{equation}\frac{dy}{dx} = f(x)\varphi(y) \tag{2.1}\end{equation}
的方程称为\textbf{变量分离方程}，其中 $f(x), \varphi(y)$ 分别是 $x, y$ 的连续函数。
\end{mydefinition}

\subsubsection{求解方法}

当 $\varphi(y) \neq 0$ 时，将 (2.1) 写成
\begin{equation}\frac{dy}{\varphi(y)} = f(x)\,dx \tag{2.2}\end{equation}
两边积分得
$$\int \frac{dy}{\varphi(y)} = \int f(x)\,dx + c$$
由 (2.2) 所确定的函数 $\Phi(x, y) = c$ 即为 (2.1) 的解。

\begin{remark}
若存在 $y_0$ 使 $\varphi(y_0) = 0$，则 $y = y_0$ 也是方程的解，但它不包含在通解中，\textbf{必须予以补上}。
\end{remark}

\begin{example}
求微分方程 $\dfrac{dy}{dx} = y(10 - y)$ 的所有解。
\end{example}

\begin{solution}
方程两边同除以 $y(10-y)$，分离变量得
$$\frac{dy}{y(10-y)} = dx$$
积分得
$$\ln\left|\frac{y}{10-y}\right| = x + c_1$$
从上式中解出 $y$（将常数记为 $c \neq 0$）：
$$y = \frac{10}{1 + ce^{-x}}$$
此外由 $\varphi(y) = y(10-y) = 0$ 得 $y = 0$ 和 $y = 10$ 也是方程的解。

故方程的所有解为：$y = \dfrac{10}{1+ce^{-x}}$（$c$ 为任常数）以及 $y = 0$、$y = 10$。
\end{solution}

\begin{example}
求微分方程 $\dfrac{dy}{dx} = \dfrac{2\sqrt{3}\,y}{x}$ 的通解。
\end{example}

\begin{solution}
分离变量后得 $\dfrac{dy}{y} = \dfrac{2\sqrt{3}}{x}\,dx$，两边积分得
$$\ln|y| = 2\sqrt{3}\ln|x| + c_1$$
整理后得通解 $y = Cx^{2\sqrt{3}}$（$C$ 为任常数，含 $C = 0$ 对应的零解 $y \equiv 0$）。
\end{solution}

\begin{example}
求微分方程 $\dfrac{dy}{dx} = p(x)y$ 的通解，其中 $p(x)$ 是 $x$ 的连续函数。
\end{example}

\begin{solution}
分离变量 $\dfrac{dy}{y} = p(x)\,dx$，两边积分 $\ln|y| = \int p(x)\,dx + c_1$，由对数定义
$$y = e^{c_1}e^{\int p(x)\,dx} = ce^{\int p(x)\,dx}$$
允许 $c = 0$ 时也包括 $y = 0$，故方程的通解为
$$y = ce^{\int p(x)\,dx} \quad (\text{$c$ 为任常数})$$
\end{solution}

\begin{example}
求初值问题 $\begin{cases} \dfrac{dy}{dx} = y^2\cos x \\ y(0) = 1 \end{cases}$ 的特解。
\end{example}

\begin{solution}
先求通解。当 $y \neq 0$ 时分离变量 $\dfrac{dy}{y^2} = \cos x\,dx$，两边积分得 $-\dfrac{1}{y} = \sin x + c$，通解为
$$y = -\frac{1}{\sin x + c}$$
此外 $y = 0$ 也是方程的解。代入初始条件 $y(0) = 1$ 得 $c = -1$，故所求特解为
$$y = \frac{1}{1 - \sin x}$$
\end{solution}

\subsection{可化为变量分离方程的类型}

\subsubsection{(I) 齐次方程}

\textbf{形如}
\begin{equation}\frac{dy}{dx} = g\left(\frac{y}{x}\right) \tag{2.5}\end{equation}
的方程称为\textbf{齐次方程}，其中 $g(u)$ 是 $u$ 的连续函数。

\textbf{求解方法}：作变量代换 $u = \dfrac{y}{x}$（即 $y = xu$），由 $y' = u + x\dfrac{du}{dx}$ 得
$$x\frac{du}{dx} = g(u) - u$$
这是关于 $u$ 的变量分离方程，解出后再变量还原。

\begin{example}
求解方程 $x\dfrac{dy}{dx} = y + \sqrt{xy}$（$x < 0$）。
\end{example}

\begin{solution}
方程变形为 $\dfrac{dy}{dx} = \dfrac{y}{x} + \sqrt{\dfrac{y}{x}}$，这是齐次方程。令 $u = \dfrac{y}{x}$，代入得 $x\dfrac{du}{dx} = \sqrt{u}$，分离变量 $\dfrac{du}{\sqrt{u}} = \dfrac{dx}{x}$，积分得
$$2\sqrt{u} = \ln|x| + c$$
变量还原得原方程的通解 $2\sqrt{\dfrac{y}{x}} = \ln|x| + c$（$x<0$）。
\end{solution}

\begin{example}
求初值问题 $\begin{cases} (y^2 + x^2)\,dx + xy\,dy = 0 \\ y(1) = 0 \end{cases}$ 的解。
\end{example}

\begin{solution}
方程变形为 $\dfrac{dy}{dx} = -\dfrac{x^2+y^2}{xy} = -\dfrac{1+(y/x)^2}{y/x}$，这是齐次方程。令 $u = \dfrac{y}{x}$，代入得 $x\dfrac{du}{dx} = -\dfrac{1+u^2}{u} - u = -\dfrac{1+2u^2}{u}$，分离变量并积分：
$$\int \frac{u}{1+2u^2}\,du = -\int \frac{dx}{x}$$
得 $\dfrac{1}{4}\ln(1+2u^2) = -\ln|x| + c$，整理得 $x^4(1+2u^2) = C$，即 $x^4 + 2x^2y^2 = C$。由 $y(1)=0$ 定出 $C=1$，故初值问题的解为 $x^4 + 2x^2y^2 = 1$。
\end{solution}

\subsubsection{(II) 分式线性型方程}

其中 $a_i, b_i, c_i$ 均为常数，可经过变量变换化为变量分离方程。分三种情况：

\textbf{情形 1}：$c_1 = c_2 = 0$，方程为齐次方程
$$\frac{dy}{dx} = g\left(\frac{y}{x}\right)$$
由 (I) 可化为变量分离方程。

\textbf{情形 2}：$\dfrac{a_1}{a_2} = \dfrac{b_1}{b_2} = k$（即 $\begin{vmatrix} a_1 & b_1 \\ a_2 & b_2 \end{vmatrix} = 0$），令 $u = a_2 x + b_2 y$，则方程化为变量分离方程。

\textbf{情形 3}：$c_1, c_2$ 不同时为零，且 $\begin{vmatrix} a_1 & b_1 \\ a_2 & b_2 \end{vmatrix} \neq 0$。解方程组
$$\begin{cases} a_1x + b_1y + c_1 = 0 \\ a_2x + b_2y + c_2 = 0 \end{cases}$$
得交点 $(\alpha, \beta) \neq (0,0)$，作坐标变换
$$X = x - \alpha, \quad Y = y - \beta$$
则方程化为 $\dfrac{dY}{dX} = g\left(\dfrac{Y}{X}\right)$ 的齐次方程，进而求解。

\textbf{求解步骤}：
\begin{enumerate}
    \item 解方程组得 $x = \alpha, y = \beta$
    \item 作变换 $X = x - \alpha, Y = y - \beta$，方程化为 $\dfrac{dY}{dX} = g\left(\dfrac{Y}{X}\right)$
    \item 再经变换 $u = \dfrac{Y}{X}$ 化为变量分离方程
    \item 求解
    \item 变量还原
\end{enumerate}

\begin{example}
求微分方程 $\dfrac{dy}{dx} = \dfrac{x - y + 1}{x + y - 3}$ 的通解。
\end{example}

\begin{solution}
解方程组 $\begin{cases} x - y + 1 = 0 \\ x + y - 3 = 0 \end{cases}$ 得 $x = 1, y = 2$。令 $X = x - 1, Y = y - 2$，代入得
$$\frac{dY}{dX} = \frac{X - Y}{X + Y}$$
令 $u = \dfrac{Y}{X}$，由 $u + X\dfrac{du}{dX} = \dfrac{1-u}{1+u}$ 得分离变量式 $\dfrac{1+u}{1-2u-u^2}\,du = \dfrac{dX}{X}$，积分得
$$-\frac{1}{2}\ln|1-2u-u^2| = \ln|X| + c$$
整理得 $X^2 - 2XY - Y^2 = C$，变量还原得原方程的通解：
$$(x-1)^2 - 2(x-1)(y-2) - (y-2)^2 = C$$
\end{solution}

\textbf{更一般的类型}（可经适当变量变换化为变量分离方程）：
\begin{itemize}
    \item $\dfrac{dy}{dx} = f(ax + by + c)$ —— 令 $u = ax+by+c$
    \item $x\dfrac{dy}{dx} = f(xy)$ —— 令 $u = xy$
    \item $x\dfrac{dy}{dx} = f\left(\dfrac{y}{x^2}\right)$ —— 令 $u = \dfrac{y}{x^2}$
    \item $M(x,y)(x\,dx + y\,dy) + N(x,y)(x\,dy - y\,dx) = 0$，其中 $M, N$ 为齐次函数（次数可以不相同）
\end{itemize}

\begin{example}
求微分方程 $(xy + y^2)\,dx + (x^2 - xy)\,dy = 0$ 的通解。
\end{example}

\begin{solution}
这是齐次方程，令 $y = vx$，则 $dy = v\,dx + x\,dv$，代入得
$$x^2\left[v(1+v) + v(1-v)\right]dx + x^3(1-v)dv = 0$$
即 $2v\,dx + x(1-v)dv = 0$，分离变量并积分：
$$\frac{dx}{x} + \frac{1-v}{2v}\,dv = 0 \quad \Rightarrow \quad \ln|x| + \frac{1}{2}\ln|v| - \frac{v}{2} = c$$
整理得 $\ln|xy| - \dfrac{y}{x} = C$，即原方程的通解
$$\ln|xy| - \frac{y}{x} = C$$
\end{solution}

\textbf{作业（2.1）}：p25：1.（2、6、7）；2.（1、4、6）；4；5；7

\section{线性方程与常数变易法}

\subsection{一阶线性微分方程}

一阶线性微分方程在 $a(x) \neq 0$ 的区间上可写成
\begin{equation}\frac{dy}{dx} + P(x)y = Q(x) \tag{1}\end{equation}
其中 $P(x), Q(x)$ 在考虑的区间上是连续函数。

\begin{itemize}
    \item 若 $Q(x) \equiv 0$，得 $y' + P(x)y = 0$，称为\textbf{一阶齐次线性方程}
    \item 若 $Q(x) \not\equiv 0$，称为\textbf{一阶非齐次线性方程}
\end{itemize}

\subsection{常数变易法}

\textbf{步骤}：
\begin{enumerate}
    \item 解对应的齐次方程 $y' + P(x)y = 0$，分离变量得
    $$y = ce^{-\int P(x)\,dx}$$
    \item 常数变易：将常数 $c$ 变为待定函数 $c(x)$，令
    $$y = c(x)e^{-\int P(x)\,dx}$$
    \item 代入原方程 (1)，得 $c'(x) = Q(x)e^{\int P(x)\,dx}$，积分得
    $$c(x) = \int Q(x)e^{\int P(x)\,dx}\,dx + \tilde{c}$$
\end{enumerate}
故 (1) 的通解为
\begin{equation}y = e^{-\int P(x)\,dx}\left(\int Q(x)e^{\int P(x)\,dx}\,dx + \tilde{c}\right) \tag{3}\end{equation}

\begin{remark}
求 (1) 的通解可直接用公式 (3)。
\end{remark}

\begin{example}
求方程 $(1+x)\dfrac{dy}{dx} - ny = e^x(1+x)^{n+1}$ 的通解（$n$ 为常数）。
\end{example}

\begin{solution}
改写为 $\dfrac{dy}{dx} - \dfrac{n}{1+x}y = e^x(1+x)^n$。齐次方程通解 $y = c(1+x)^n$。常数变易，令 $y = c(x)(1+x)^n$，代入得 $c'(x) = e^x$，积分得 $c(x) = e^x + \tilde{c}$，故通解为
$$y = (1+x)^n(e^x + \tilde{c})$$
\end{solution}

\begin{example}
求方程 $\dfrac{dy}{dx} = x^2 - \dfrac{y}{x}$ 的通解。
\end{example}

\begin{solution}
改写为 $\dfrac{dy}{dx} + \dfrac{1}{x}y = x^2$，用公式 (3)（$P = \dfrac{1}{x}, Q = x^2$）：
$$y = e^{-\ln x}\left(\int x^2 e^{\ln x}\,dx + c\right) = \frac{1}{x}\left(\frac{x^4}{4} + c\right)$$
\end{solution}

\begin{example}
求初值问题 $\begin{cases} \dfrac{dy}{dx} = \dfrac{4y}{x^2} + x^2 + 1 \\ y(1) = 1 \end{cases}$ 的解。
\end{example}

\begin{solution}
原方程即 $y' - \dfrac{4}{x^2}y = x^2+1$，用公式 (3) 得通解
$$y = e^{4/x}\left(\int (x^2+1)e^{-4/x}\,dx + \tilde{c}\right)$$
其中积分 $\int(x^2+1)e^{-4/x}dx$ 无初等原函数（含指数积分），通解即保留此积分形式；代入 $y(1) = 1$ 定出常数即可。
\end{solution}

\textbf{练习}：
\begin{enumerate}
    \item $(1+x^2)y' + xy = e^{-x}$ $\rightarrow$ 通解 $y = \dfrac{1}{\sqrt{1+x^2}}\left(\int \dfrac{e^{-x}}{\sqrt{1+x^2}}\,dx + c\right)$（积分无初等原函数，通解保留积分形式）
    \item $y' + 2xy = 4x$ $\rightarrow$ 通解 $y = 2 + ce^{-x^2}$
    \item $\dfrac{dy}{dx} + y = xy^3$ $\rightarrow$ 令 $z = y^{-2}$ 化为 $z' - 2z = -2x$，通解 $y^{-2} = x + \frac{1}{2} + ce^{2x}$
\end{enumerate}

\subsection{伯努利（Bernoulli）方程}

\textbf{形如}
$$\frac{dy}{dx} + P(x)y = Q(x)y^n \quad (n \neq 0, 1)$$
的方程称为伯努利方程，其中 $P(x), Q(x)$ 为 $x$ 的连续函数。

\textbf{解法}：引入变量变换 $z = y^{1-n}$，方程化为
$$\frac{dz}{dx} + (1-n)P(x)z = (1-n)Q(x)$$
求上述线性方程的通解，再变量还原。

\begin{example}
求方程 $\dfrac{dy}{dx} + \dfrac{1}{x}y = 2y^2$ 的通解。
\end{example}

\begin{solution}
这是 $n=2$ 的伯努利方程。令 $z = y^{-1}$，方程化为 $\dfrac{dz}{dx} - \dfrac{1}{x}z = -2$，解得 $z = cx + 2$，故通解为
$$y = \frac{1}{cx + 2}$$
\end{solution}

\begin{example}
求方程 $2xy' + y = y^2\ln x$ 的通解。
\end{example}

\begin{solution}
改写为 $y' + \dfrac{1}{2x}y = \dfrac{\ln x}{2x}y^2$，这是 $n=2$ 的伯努利方程。令 $z = y^{-1}$，得 $z' - \dfrac{1}{2x}z = -\dfrac{\ln x}{2x}$，解得
$$z = \ln x + 2 + c\sqrt{x}$$
故通解 $y = \dfrac{1}{\ln x + 2 + c\sqrt{x}}$。
\end{solution}

\subsection{线性微分方程的应用举例：R-L 串联电路}

由电感 $L$、电阻 $R$ 和电源组成的串联电路，电动势 $E$ 为常数。开关 $K$ 合上后（$t=0$ 时刻，初始电流 $I(0) = 0$），求电流强度 $I(t)$。

由 Kirchhoff 第二定律（闭合回路电压代数和为零）：电感压降 $L\dfrac{dI}{dt}$，电阻压降 $RI$，故
$$L\frac{dI}{dt} + RI = E \quad \Rightarrow \quad \frac{dI}{dt} + \frac{R}{L}I = \frac{E}{L}$$
解线性方程得通解 $I(t) = \dfrac{E}{R} + ce^{-Rt/L}$，由 $I(0) = 0$ 得 $c = -\dfrac{E}{R}$，故
$$I(t) = \frac{E}{R}\left(1 - e^{-Rt/L}\right)$$

\textbf{作业（2.2）}：p30 1.（2、4、8、12）；2；p31 4；5.（2）；7.（2）

\section{恰当方程与积分因子}

\subsection{恰当方程的定义及条件}

设 $u(x,y)$ 是连续可微的函数，其全微分为 $du = \dfrac{\partial u}{\partial x}dx + \dfrac{\partial u}{\partial y}dy$。如果恰好有方程
$$\frac{\partial u(x,y)}{\partial x}dx + \frac{\partial u(x,y)}{\partial y}dy = 0$$
就可以马上写出它的隐式解 $u(x,y) = c$。

\begin{mydefinition}{定义}
若有函数 $u(x,y)$ 使得 $du(x,y) = M(x,y)dx + N(x,y)dy$，则称微分方程
\begin{equation}M(x,y)dx + N(x,y)dy = 0 \tag{1}\end{equation}
是\textbf{恰当方程}。此时 (1) 的通解为 $u(x,y) = c$。
\end{mydefinition}

\begin{example}
$xdy + ydx = 0$ 是恰当方程（$= d(xy)$）；$(2xy)dx + (x^2 + 3y^2)dy = 0$ 是恰当方程（$= d(x^2y + y^3)$）。
\end{example}

\textbf{需考虑的问题}：
\begin{enumerate}
    \item 方程 (1) 是否为恰当方程？
    \item 若是恰当方程，怎样求解？
    \item 若不是恰当方程，有无可能转化为恰当方程求解？
\end{enumerate}

\subsection{方程为恰当方程的充要条件}

\begin{myTheorem}{定理}
设函数 $M(x,y), N(x,y)$ 在矩形区域 $R$ 内连续且有连续的一阶偏导数，则方程 $M(x,y)dx + N(x,y)dy = 0$ 为恰当方程的充要条件是
\begin{equation}\frac{\partial M(x,y)}{\partial y} = \frac{\partial N(x,y)}{\partial x} \tag{2}\end{equation}
\end{myTheorem}

\begin{proof}
\textbf{必要性}：设 (1) 是恰当方程，存在 $u(x,y)$ 使 $du = Mdx + Ndy$，则 $\dfrac{\partial u}{\partial x} = M, \dfrac{\partial u}{\partial y} = N$，从而 $\dfrac{\partial^2 u}{\partial y\partial x} = \dfrac{\partial M}{\partial y}$，$\dfrac{\partial^2 u}{\partial x\partial y} = \dfrac{\partial N}{\partial x}$，由混合偏导数连续相等得 $\dfrac{\partial M}{\partial y} = \dfrac{\partial N}{\partial x}$。

\textbf{充分性}：若 $\dfrac{\partial M}{\partial y} = \dfrac{\partial N}{\partial x}$，构造 $u(x,y) = \int M(x,y)\,dx + \varphi(y)$，其中 $\varphi(y)$ 由 $\dfrac{\partial u}{\partial y} = N$ 确定：$\varphi'(y) = N - \dfrac{\partial}{\partial y}\int M\,dx$，可证右端与 $x$ 无关，积分即得。
\end{proof}

\textbf{通解公式}：若 (1) 为恰当方程，则其通解为
$$\int M(x,y)\,dx + \int\left(N - \frac{\partial}{\partial y}\int M\,dx\right)dy = c$$

\subsection{恰当方程的求解}

关键是求原函数 $u(x,y)$ 的问题。三种方法：

\subsubsection{1. 不定积分法}

步骤：\textcircled{1} 判断是否恰当方程；\textcircled{2} 设 $u(x,y) = \int M\,dx + \varphi(y)$；\textcircled{3} 由 $\dfrac{\partial u}{\partial y} = N$ 求 $\varphi(y)$。

\begin{example}
验证方程 $(e^x + y)\,dx + (x - 2\sin y)\,dy = 0$ 是恰当方程并求通解。
\end{example}

\begin{solution}
$M = e^x + y, N = x - 2\sin y$，$\dfrac{\partial M}{\partial y} = 1 = \dfrac{\partial N}{\partial x}$，故是恰当方程。

设 $u = \int(e^x + y)dx = e^x + yx + \varphi(y)$，由 $\dfrac{\partial u}{\partial y} = x + \varphi'(y) = x - 2\sin y$ 得 $\varphi'(y) = -2\sin y$，积分 $\varphi(y) = 2\cos y$。故通解
$$e^x + xy + 2\cos y = c$$
\end{solution}

\subsubsection{2. 线积分法}

由曲线积分与路径无关定理，取 $(x_0, y_0) \in R$：
$$u(x,y) = \int_{x_0}^{x} M(x, y_0)\,dx + \int_{y_0}^{y} N(x, y)\,dy$$
通解为 $\int_{x_0}^{x} M(x,y_0)\,dx + \int_{y_0}^{y} N(x,y)\,dy = c$。

\begin{example}
求解方程 $(2x\cos y + e^y)\,dx + (e^y - x^2\sin y)\,dy = 0$。
\end{example}

\begin{solution}
$M = 2x\cos y + e^y, N = e^y - x^2\sin y$，$\dfrac{\partial M}{\partial y} = -2x\sin y + e^y = \dfrac{\partial N}{\partial x}$，是恰当方程。取 $(x_0,y_0) = (0,0)$：
$$u = \int_0^x 2x\,dx + \int_0^y (e^y - x^2\sin y)\,dy = x^2 + e^y - 1 + x^2\cos y - x^2$$
整理得通解
$$x^2\cos y + e^y = c$$
\end{solution}

\subsubsection{3. 分组凑微法}

采用"分项组合"的方法，把本身已构成全微分的项分出来，再把剩余项凑成全微分。应熟记简单二元函数的全微分：
$$xdy + ydx = d(xy), \quad \frac{xdy - ydx}{y^2} = d\left(\frac{x}{y}\right), \quad \frac{ydx - xdy}{x^2} = d\left(\frac{y}{x}\right)$$
$$\frac{xdy - ydx}{x^2 + y^2} = d\left(\arctan\frac{y}{x}\right), \quad \frac{xdy - ydx}{xy} = d\left(\ln\left|\frac{y}{x}\right|\right), \quad \frac{xdy - ydx}{x^2 - y^2} = \frac{1}{2}d\left(\ln\left|\frac{x+y}{x-y}\right|\right)$$

\begin{example}
求方程 $(3x^2 + 6xy^2)\,dx + (6x^2y + 4y^3)\,dy = 0$ 的通解。
\end{example}

\begin{solution}
$\dfrac{\partial M}{\partial y} = 12xy = \dfrac{\partial N}{\partial x}$，是恰当方程。分项组合：
$$(3x^2\,dx + 4y^3\,dy) + (6xy^2\,dx + 6x^2y\,dy) = d(x^3) + d(y^4) + d(3x^2y^2) = 0$$
故通解为 $x^3 + y^4 + 3x^2y^2 = c$。
\end{solution}

\begin{example}
验证方程 $(\cos x - x\sin x + xy^2)\,dx + (x^2y - y)\,dy = 0$ 是恰当方程，并求满足 $y(0) = 2$ 的解。
\end{example}

\begin{solution}
$M = \cos x - x\sin x + xy^2$，$N = x^2y - y$，$\dfrac{\partial M}{\partial y} = 2xy = \dfrac{\partial N}{\partial x}$，是恰当方程。分项组合得 $d(x\cos x) + d\left(\dfrac{x^2y^2}{2}\right) - d\left(\dfrac{y^2}{2}\right) = 0$，通解为
$$x\cos x + \frac{x^2y^2}{2} - \frac{y^2}{2} = c$$
由 $y(0) = 2$ 得 $c = -2$，故 $x\cos x + \dfrac{x^2y^2 - y^2}{2} = -2$，即 $x^2y^2 - y^2 + 2x\cos x + 4 = 0$。
\end{solution}

\subsection{积分因子}

对非恰当方程，乘上一个因子后可变为恰当方程。

\begin{example}
一阶线性方程 $dy - (P(x)y + Q(x))dx = 0$ 不是恰当方程，乘以 $e^{-\int P(x)dx}$ 后变为恰当方程。
\end{example}

\subsubsection{1. 定义}

\begin{mydefinition}{定义}
如果存在连续可微函数 $\mu(x,y) \neq 0$，使得
$$\mu(x,y)M(x,y)\,dx + \mu(x,y)N(x,y)\,dy = 0$$
为恰当方程，则称 $\mu(x,y)$ 是方程 (1) 的一个\textbf{积分因子}。
\end{mydefinition}

\begin{example}
验证 $\mu = xy$ 是方程 $(3xy + 4y^2)\,dx + (2x^2 + 3xy)\,dy = 0$ 的积分因子并求通解。
\end{example}

\begin{solution}
乘以 $\mu = xy$ 后 $M' = 3x^2y^2 + 4xy^3, N' = 2x^3y + 3x^2y^2$，$\dfrac{\partial M'}{\partial y} = 6x^2y + 12xy^2 = \dfrac{\partial N'}{\partial x}$，是恰当方程。分项组合得 $d(x^3y^2) + d(x^2y^3) = 0$，通解
$$x^3y^2 + x^2y^3 = c$$
\end{solution}

\subsubsection{2. 积分因子的确定}

$\mu(x,y)$ 是方程 (1) 的积分因子的充要条件是
$$\mu\left(\frac{\partial M}{\partial y} - \frac{\partial N}{\partial x}\right) = N\frac{\partial \mu}{\partial x} - M\frac{\partial \mu}{\partial y}$$
这本身是偏微分方程，一般比原方程更困难，但提供了寻找特殊形式积分因子的途径。

\textbf{仅依赖于 $x$ 的积分因子}：方程有仅依赖于 $x$ 的积分因子 $\mu = \mu(x)$ 的充要条件是
$$\psi(x) = \frac{\dfrac{\partial M}{\partial y} - \dfrac{\partial N}{\partial x}}{N}$$
只是 $x$ 的函数（与 $y$ 无关），此时积分因子为
\begin{equation}\mu(x) = e^{\int \psi(x)\,dx} \tag{10}\end{equation}

\textbf{仅依赖于 $y$ 的积分因子}：同理，充要条件是
$$\varphi(y) = \frac{\dfrac{\partial N}{\partial x} - \dfrac{\partial M}{\partial y}}{M}$$
只是 $y$ 的函数，此时积分因子为 $\mu(y) = e^{\int \varphi(y)\,dy}$。

\begin{example}
求方程 $(2y + 2xy + 2y^2)\,dx + (2x + 4y)\,dy = 0$ 的通解。
\end{example}

\begin{solution}
$M = 2y + 2xy + 2y^2, N = 2x + 4y$，$\dfrac{\partial M}{\partial y} = 2 + 2x + 4y$，$\dfrac{\partial N}{\partial x} = 2$，不是恰当方程。计算
$$\psi(x) = \frac{\dfrac{\partial M}{\partial y} - \dfrac{\partial N}{\partial x}}{N} = \frac{2x + 4y}{2x + 4y} = 1$$
仅依赖 $x$，故积分因子为 $\mu = e^{\int 1\,dx} = e^x$。乘以 $\mu$ 后
$$e^x(2y + 2xy + 2y^2)dx + e^x(2x + 4y)dy = 0$$
分项组合得 $d(2xye^x) + d(2y^2e^x) = 0$，通解为
$$2xye^x + 2y^2e^x = c$$
\end{solution}

\begin{example}
求解方程 $\dfrac{dy}{dx} = -\dfrac{x + \sqrt{x^2 + y^2}}{y}$（$y > 0$）。
\end{example}

\begin{solution}
改写为 $x\,dx + y\,dy + \sqrt{x^2+y^2}\,dx = 0$，即 $\dfrac{1}{2}d(x^2+y^2) + \sqrt{x^2+y^2}\,dx = 0$，取积分因子 $\mu = \dfrac{1}{\sqrt{x^2+y^2}}$，得
$$d\left(\sqrt{x^2+y^2}\right) + dx = 0$$
故通解为 $\sqrt{x^2+y^2} + x = c$。
\end{solution}

\begin{example}
求解方程 $ydx - (x + y)\,dy = 0$ 的通解（四种方法）。
\end{example}

\begin{solution}
\begin{itemize}
    \item \textbf{方法 1（积分因子）}：$\dfrac{\partial M}{\partial y} - \dfrac{\partial N}{\partial x} = 1 - (-1) = 2$，$\dfrac{\dfrac{\partial N}{\partial x} - \dfrac{\partial M}{\partial y}}{M} = -\dfrac{2}{y}$ 仅依赖 $y$，$\mu = e^{\int -2/y\,dy} = \dfrac{1}{y^2}$。乘以 $\mu$ 得 $\dfrac{y\,dx - x\,dy}{y^2} - \dfrac{dy}{y} = 0$，即 $d\left(\dfrac{x}{y}\right) - d(\ln y) = 0$，通解 $\dfrac{x}{y} - \ln y = c$，即 $x = y\ln y + cy$。
    \item \textbf{方法 2（凑微分）}：$ydx - xdy = y\,dy$，$\dfrac{ydx - xdy}{y^2} = \dfrac{dy}{y}$，得 $d\left(\dfrac{x}{y}\right) = d(\ln y)$，通解 $x = y\ln y + cy$。
    \item \textbf{方法 3（齐次方程）}：$\dfrac{dy}{dx} = \dfrac{y}{x + y}$ 是齐次方程，令 $u = \dfrac{y}{x}$，由 $u + xu' = \dfrac{u}{1+u}$ 得 $xu' = -\dfrac{u^2}{1+u}$，分离变量 $\dfrac{1+u}{u^2}\,du = -\dfrac{dx}{x}$，积分得 $\ln|u| - \dfrac{1}{u} = -\ln|x| + c$，即 $\ln|y| - \dfrac{x}{y} = c$，变量还原得 $x = y\ln y + cy$。
    \item \textbf{方法 4（线性方程）}：改写为 $\dfrac{dx}{dy} - \dfrac{1}{y}x = 1$，以 $y$ 为自变量、$x$ 为未知函数的线性方程，公式 (3) 得 $x = y\ln y + cy$。
\end{itemize}
\end{solution}

\textbf{作业（2.3）}：p41 1.（1）（5）；2.（1）（4）（7）（9）；3、4、5、6、9

\section{一阶隐方程与参数表示}

一阶隐式方程 $F(x, y, y') = 0$，当 $y'$ 未能解出或解出形式复杂时，采用引进参数的办法使其变为导数已解出的方程类型。主要研究四种类型：
\begin{enumerate}
    \item $y = f(x, y')$
    \item $x = f(y, y')$
    \item $F(x, y') = 0$（不显含 $y$）
    \item $F(y, y') = 0$（不显含 $x$）
\end{enumerate}

\textbf{参数形式解的定义}：如果存在定义在 $(\alpha, \beta)$ 上的函数 $\varphi(t), \psi(t)$ 使 $x = \varphi(t), y = \psi(t)$，且 $F(\varphi(t), \psi(t), \psi'(t)/\varphi'(t)) = 0$，则称其为方程的参数形式解。

\subsection{可解出 $y'$ 的方程：$y = f(x, y')$}

\noindent 设 $y' = p$，方程变为 $y = f(x, p)$，两边对 $x$ 求导，以 $p' = \dfrac{dp}{dx}$ 代入：
$$p = \frac{\partial f}{\partial x} + \frac{\partial f}{\partial p}\frac{dp}{dx}$$
这是关于变量 $x, p$ 的一阶微分方程。求出其通解后回代：
\begin{itemize}
    \item 若通解为 $p = \varphi(x, c)$，代入 $y = f(x, p)$ 即得通解 $y = f(x, \varphi(x,c))$
    \item 若通解为 $x = \psi(p, c)$，得参数形式通解 $\begin{cases} x = \psi(p, c) \\ y = f(\psi(p,c), p) \end{cases}$
    \item 若通解为 $\Phi(x, p, c) = 0$，得参数形式通解 $\begin{cases} \Phi(x, p, c) = 0 \\ y = f(x, p) \end{cases}$
\end{itemize}

\begin{example}
求解方程 $y = x\dfrac{dy}{dx} + \left(\dfrac{dy}{dx}\right)^2$ 的通解。
\end{example}

\begin{solution}
令 $p = y'$，则 $y = xp + p^2$，两边对 $x$ 求导：
$$p = p + x\frac{dp}{dx} + 2p\frac{dp}{dx} = p + (x + 2p)\frac{dp}{dx}$$
即 $(x + 2p)\dfrac{dp}{dx} = 0$。

\begin{itemize}
    \item 由 $\dfrac{dp}{dx} = 0$ 得 $p = c$，代入得通解 $y = cx + c^2$
    \item 由 $x + 2p = 0$ 得 $p = -\dfrac{x}{2}$，代入得 $y = -\dfrac{x^2}{4}$
\end{itemize}
曲线 $y = -\dfrac{x^2}{4}$ 是积分曲线族的\textbf{包络}，在微分方程中称 $y = -\dfrac{x^2}{4}$ 为原方程的\textbf{奇解}。

\begin{remark}
讲义原题 $y = p^2 - 2xp + 2x^2$ 与其通解 $y = c^2 + 2cx + x^2$、奇解 $y = 2x^2$ 相互矛盾（通解代入方程右端得 $4c^2+4cx+2x^2 \neq c^2+2cx+x^2$，奇解代入亦不成立）。按讲义通解 $y = (x+c)^2$ 反推，正确原题应为 Clairaut 方程 $y = xp + p^2$，此处按此修正。
\end{remark}
\end{solution}

\subsection{可解出 $x$ 的方程：$x = f(y, y')$}

设 $y' = p$，方程变为 $x = f(y, p)$，两边对 $y$ 求导，以 $\dfrac{dx}{dy} = \dfrac{1}{p}$ 代入：
$$\frac{1}{p} = \frac{\partial f}{\partial y} + \frac{\partial f}{\partial p}\frac{dp}{dy}$$
得到关于变量 $y, p$ 的一阶微分方程，求出通解 $\Phi(y, p, c) = 0$ 后得参数形式通解 $\begin{cases} \Phi(y,p,c) = 0 \\ x = f(y, p) \end{cases}$。

\begin{example}
求解方程 $x\left(\dfrac{dy}{dx}\right)^3 - y\dfrac{dy}{dx} + 2 = 0$。
\end{example}

\begin{solution}
设 $p = y'$，则 $x = \dfrac{yp - 2}{p^3}$（$p \neq 0$），即 $y = xp^2 + \dfrac{2}{p}$，两边对 $y$ 求导（$\dfrac{dx}{dy} = \dfrac{1}{p}$）：
$$1 = p^2\frac{dx}{dy} + \left(2xp - \frac{2}{p^2}\right)\frac{dp}{dy}$$
由 $\dfrac{dp}{dy} = \dfrac{1-p}{2xp - 2/p^2}$ 得关于 $x, p$ 的一阶线性方程
$$x' - \frac{2}{1-p}x = -\frac{2}{p^3(1-p)}$$
积分因子为 $(1-p)^2$，解得
$$x = \frac{1}{(1-p)^2}\left(\frac{1}{p^2} - \frac{2}{p} + C\right), \qquad y = xp^2 + \frac{2}{p}$$
另有解 $y = 0$。
\end{solution}

\subsection{不显含 $y$ 的方程：$F(x, y') = 0$}

设 $y' = p$，方程 $F(x, p) = 0$ 表示 $xp$ 平面上的一条曲线。若能找到参数表示 $x = \varphi(t), p = \psi(t)$（满足 $F(\varphi(t), \psi(t)) = 0$），则由 $dy = p\,dx = \psi(t)\varphi'(t)\,dt$ 积分得
$$y = \int \psi(t)\varphi'(t)\,dt + c$$
参数形式通解为 $\begin{cases} x = \varphi(t) \\ y = \int \psi(t)\varphi'(t)\,dt + c \end{cases}$。关键一步（也是最困难一步）是找曲线的参数表示。

\begin{example}
求解方程 $y' = x\sqrt{1 + (y')^2}$。
\end{example}

\begin{solution}
设 $p = y'$，则 $p = x\sqrt{1+p^2}$，即 $x = \dfrac{p}{\sqrt{1+p^2}}$。取参数 $t$：令 $p = \tan t$，则 $x = \sin t$。由 $dy = p\,dx = \tan t \cos t\,dt = \sin t\,dt$，积分得 $y = -\cos t + c$。参数形式通解
$$\begin{cases} x = \sin t \\ y = c - \cos t \end{cases}$$
消去参数得 $x^2 + (y - c)^2 = 1$，即单位圆族。
\end{solution}

\subsection{不显含 $x$ 的方程：$F(y, y') = 0$}

设 $y' = p$，方程 $F(y, p) = 0$ 表示 $yp$ 平面上的一条曲线，参数表示 $y = \varphi(t), p = \psi(t)$。由 $dx = \dfrac{dy}{p} = \dfrac{\varphi'(t)}{\psi(t)}\,dt$ 积分得
$$x = \int \frac{\varphi'(t)}{\psi(t)}\,dt + c$$
参数形式通解为 $\begin{cases} y = \varphi(t) \\ x = \int \dfrac{\varphi'(t)}{\psi(t)}\,dt + c \end{cases}$。

\begin{example}
求解方程 $y = (y')^2 - 2y' + 1$。
\end{example}

\begin{solution}
设 $p = y'$，方程即 $y = p^2 - 2p + 1 = (p-1)^2$，取 $p = t$ 为参数，则 $y = (t-1)^2$。由 $dx = \dfrac{dy}{p} = \dfrac{2(t-1)}{t}\,dt = 2\left(1 - \dfrac{1}{t}\right)dt$，积分 $x = 2t - 2\ln|t| + c$。参数形式通解
$$\begin{cases} y = (t-1)^2 \\ x = 2t - 2\ln|t| + c \end{cases}$$
此外 $y = 1$ 也是方程的解（$y' = 0$ 时右端 $= 1 = y$）。
\end{solution}

\subsection{利用变量代换的微分方程积分法}

当 $F(x, y, y') = 0$ 中 $x, y, y'$ 都不易解出，或虽能解出但积分较复杂时，可先作适当的变量代换再求解。如何选择变量代换没有一定规律，需要大量练习积累经验。

\begin{example}
求解方程 $\dfrac{dy}{dx} = (x + y)^2$ 的通解。
\end{example}

\begin{solution}
令 $u = x + y$，则 $u' = 1 + y' = 1 + u^2$，分离变量并积分：
$$\frac{du}{1 + u^2} = dx \quad \Rightarrow \quad \arctan u = x + c$$
变量还原得通解
$$y = \tan(x + c) - x$$
\end{solution}

\begin{example}
求方程 $y(y')^2 - 2xy' + y = 0$ 的通解（令 $u = y^2, v = x^2$）。
\end{example}

\begin{solution}
令 $u = y^2, v = x^2$，则 $du = 2y\,dy, dv = 2x\,dx$，$y' = \dfrac{x}{y}\dfrac{du}{dv}$。代入原方程并整理得关于 $u, v$ 的方程
$$u = 2v\frac{du}{dv} - v\left(\frac{du}{dv}\right)^2$$
令 $p = \dfrac{du}{dv}$，则 $u = v(2p - p^2)$，两边对 $v$ 求导：
$$p = (2p - p^2) + v(2 - 2p)p' \quad \Rightarrow \quad (1 - p)(p + 2vp') = 0$$
\begin{itemize}
    \item 由 $p = 1$ 得 $u = v$，即 $y^2 = x^2$；
    \item 由 $p + 2vp' = 0$ 分离变量得 $p = \dfrac{C}{\sqrt{v}}$，代入 $u = v(2p-p^2)$ 得
    $$u = 2C\sqrt{v} - C^2, \qquad \text{即} \quad y^2 = 2Cx - C^2$$
\end{itemize}
故通解为 $y^2 = 2Cx - C^2$（$C$ 为任常数），曲线族 $y^2 = 2Cx - C^2$ 对参数 $C$ 求导得 $C = x$，代入得包络（奇解）$y^2 = x^2$，即 $y = \pm x$。
\end{solution}

\textbf{作业（2.4）}：p49-50：1.（15、19、20、25、29、30）；2、3、4；5（只做证明部分）

\section*{本章小结}

\begin{itemize}
    \item 一阶方程初等解法体系：\textbf{变量分离 $\rightarrow$ 齐次方程 $\rightarrow$ 线性方程（常数变易法）$\rightarrow$ 伯努利方程 $\rightarrow$ 恰当方程（积分因子）$\rightarrow$ 隐方程（参数法）}
    \item 能初等解法的微分方程很有限，如 \textbf{Riccati 方程} $\dfrac{dy}{dx} = P(x)y^2 + Q(x)y + R(x)$ 一般没有初等解法
\end{itemize}
